\documentclass[11pt]{amsart}

\usepackage[T1]{fontenc}
\usepackage{lmodern}
\usepackage{microtype}
\usepackage{amsmath,amssymb}
\usepackage[hidelinks]{hyperref}

\hypersetup{
  pdftitle={Counterexamples to the Equality Clause in Conjecture 20 of Chen--Guo--Li--Wang}
}

\newtheorem{theorem}{Theorem}

\title[Counterexamples to Conjecture 20]
{Counterexamples to the Equality Clause in Conjecture~20 of Chen--Guo--Li--Wang}

\subjclass[2020]{05C50}
\keywords{Laplacian eigenvalue, graph complement, Nordhaus--Gaddum inequality}
\date{}

\begin{document}

\begin{abstract}
Chen, Guo, Li, and Wang conjectured an upper bound for the product of
corresponding Laplacian eigenvalues of a graph and its complement, together
with an equality characterization by a specified join construction. We show
that the equality characterization is false. For every integer $t\ge2$, the
graph $K_{2t,2t}$ attains equality at $n=4t$ and $k=3t$, although the proposed
equality family is empty for these parameters. In particular, $K_{4,4}$ gives
an eight-vertex counterexample. The product inequality itself is not
contradicted.
\end{abstract}

\maketitle

Let $G$ be a simple graph on $n$ vertices, write $\overline G$ for its
complement, and order its Laplacian eigenvalues as
\[
\mu_1(G)\ge\mu_2(G)\ge\cdots\ge\mu_n(G)=0.
\]
For graphs $X$ and $Y$, write $X\vee Y$ for their join. Chen, Guo, Li, and
Wang~\cite[Conjecture~20]{CGLW} conjectured that, for every integer $k$
with $1\le k\le 3n/4$,
\[
\mu_k(G)\mu_k(\overline G)\le n(n-k),
\]
with equality if and only if $G$ or $\overline G$ is isomorphic to
$K_k\vee H$, where $H$ is a disconnected graph on $n-k$ vertices with at
least $k+1$ connected components.

They proved the inequality for $n/2\le k\le3n/4$ in
\cite[Theorem~18]{CGLW}. Their argument gives
\[
\mu_k(G)\mu_k(\overline G)\le \frac{n^2}{4}\le n(n-k),
\]
so equality in this range can occur only at the endpoint $k=3n/4$. The
following family shows that the proposed characterization omits equality
cases at that endpoint.

\begin{theorem}
For every integer $t\ge2$, set
\[
n=4t,\qquad k=3t,\qquad G=K_{2t,2t}.
\]
Then
\[
\mu_k(G)\mu_k(\overline G)=n(n-k),
\]
but neither $G$ nor $\overline G$ belongs to the equality family stated in
Conjecture~20.
\end{theorem}

\begin{proof}
Put $m=2t$. Since $\overline{K_{m,m}}\cong 2K_m$, where $2K_m$ denotes
the disjoint union of two copies of $K_m$, the Laplacian spectra are
\[
\operatorname{Spec}_L(K_{m,m})
 =\{2m,m^{(2m-2)},0\},
\qquad
\operatorname{Spec}_L(2K_m)
 =\{m^{(2m-2)},0^{(2)}\},
\]
where superscripts denote multiplicities. Substituting $m=2t$ gives
\[
\operatorname{Spec}_L(G)
 =\{4t,(2t)^{(4t-2)},0\},
\qquad
\operatorname{Spec}_L(\overline G)
 =\{(2t)^{(4t-2)},0^{(2)}\}.
\]
Because $t\ge2$, we have $3t\le4t-2$, and therefore
\[
\mu_{3t}(G)=\mu_{3t}(\overline G)=2t.
\]
Consequently,
\[
\mu_k(G)\mu_k(\overline G)=4t^2=(4t)(4t-3t)=n(n-k).
\]

For these parameters, a graph $H$ in the proposed equality family would have
$n-k=t$ vertices and at least $k+1=3t+1$ connected components. This is
impossible, since a graph on $t$ vertices has at most $t$ connected
components.
\end{proof}

The refutation does not depend on how many components the clause demands. Any
graph of the form $K_j\vee H$ with $j\ge1$ has $j$ vertices of degree $n-1$,
while the maximum degrees of $K_{2t,2t}$ and of $2K_{2t}$ are $2t$ and $2t-1$,
both smaller than $n-1=4t-1$. So neither $G$ nor $\overline G$ is a join
$K_j\vee H$ for any $j\ge1$, and the equality case above remains outside the
stated family under any weaker requirement on the number of components of $H$,
such as ``at least two''.

Taking $t=2$ gives
\[
(n,k,G)=(8,6,K_{4,4}),\qquad \overline G\cong2K_4,
\]
and
\[
\mu_6(K_{4,4})\mu_6(2K_4)=4\cdot4=16=8(8-6).
\]
Thus the ``only if'' direction of the equality statement in
Conjecture~20 fails. These examples attain the proposed upper bound and do
not violate the product inequality.

Chen, Guo, Li and Wang state (\cite[Remark~1]{CGLW}, where the conjecture
carries the stale draft label ``conjecture~3'') that they verified
Conjecture~20 for all graphs on at most nine vertices. The counterexample
above has eight vertices, so that verification cannot have been correct on
the equality clause. We generated all graphs on $2\le n\le9$ vertices with
nauty's \texttt{geng} and evaluated $\mu_k(G)\mu_k(\overline G)$ for
$1\le k\le\lfloor3n/4\rfloor$. Over this range the inequality
$\mu_k(G)\mu_k(\overline G)\le n(n-k)$ is never violated, and the equality
characterization fails in exactly five cases, all at $n=8$, $k=6$. Up to complementation these are $\{K_{4,4},2K_4\}$, the
complementary pair with graph6 codes \texttt{G?zvf\_}/\texttt{GQhTUg}, and
the self-complementary graph \texttt{GCXnf\_}. In particular $K_{4,4}$ is a
counterexample of minimum order. The verification program accompanying this
note does not repeat that census: it re-enumerates the graphs of order
$2\le n\le6$ up to isomorphism in exact arithmetic, finding every equality case
there inside the conjectured family, and it checks the five graphs listed above
one by one. The exhaustiveness of that list and the minimality of $K_{4,4}$
therefore rest on the \texttt{geng} search alone.

\begin{thebibliography}{1}

\bibitem{CGLW}
Q.~Chen, J.-M.~Guo, W.-J.~Li, and Z.~Wang,
\emph{A conjecture on the Nordhaus--Gaddum product type inequality for
Laplacian eigenvalues of a graph},
Electron. J. Combin. \textbf{32} (2025), no.~4, Paper P4.37, 13 pp.,
\href{https://doi.org/10.37236/14060}{doi:10.37236/14060}.

\end{thebibliography}

\end{document}